\documentclass{article}
\usepackage{graphicx} % Required for inserting images
\usepackage{amsmath}
\usepackage{amsfonts}
\usepackage{xcolor}

\begin{document}

\color{purple}{¿Cuál es la razón por la que elegiste la carrera que estás estudiando?} \\
\color{black}{Desde la secundaria me ha interesado muchísimo los tópicos de la física, y de las ciencias en general. Si bien mi gusto por las ciencias biomédicas y bioquímicas es muy grande la física siempre ha sobresalido, mis ídolos eran QuantummFracture, Javier Santaolalla, C de Ciencia entre otros. En específico los temas sobre el espacio siempre me apasionaron demasiado, por lo que mi intención al entrar a la carrera era volverme Astrofísico Relativista y posteriormente hacer posgrados en Ciencias de la Ingeniería Aeronáutica y Espacial con toda la preparación que tendría por detrás, más que nada porque siempre he querido ser un ingeniero de los de antes (físico experimental) que no se quedaba con lo ya existente sino que siempre busca explotar al máximo las más recientes vanguardias o las áreas sin explorar.} \\
\color{purple}{¿Por qué consideras que es esencial el álgebra lineal en tu carrera?} \\
\color{black}{Es fundamental porque a partir de quinto semestre, los modelos físicos probabilísticos se construyen alrededor de vectores en el espacio de Sobolev (al menos para IFC), además para ecuaciones diferenciales ordinarias es muy misterioso saber qué está sucediendo detrás de muchas herramientas matemáticas sin conocer conceptos como linealmente independencia, wronksianos, generadores. Igualmente en variable compleja 1 y 2 los teoremas de lineal como la representación matricial de cualquier aplicación lineal, relaciones de equivalencia (fundamentales para construcción de homotopía del ambiente, ciclos homólogos entre sí o a cero), etc. Para cálculo diferencial e integral univariable y vectorial en series de Fourier (incluyendo su versión compleja), pues aplicarlo sin conocer sobre el producto interior definido por la integral que induce cada coeficiente es perderse de mucha historia detrás. Para cálculo variacional sirve en describir espacios de funciones, algunos lemas fuertes de funcionales previos a Euler-Lagrange, entre otras aplicaciones...} \\
\color{purple}{¿En qué sector esperas desarrollar tu carrera profesional?} \\
\color{black}{En la industria astronáutica.} \\
\color{purple}{¿Cómo relacionarías la pregunta anterior con la materia de computación cuántica?} \\
\color{black}{Debo decir previamente que la computación cuántica me ha supuesto una revolución en la humanidad que a pasos agigantados cambiará todo tal como lo vemos, por lo que mi nuevo tentativo interés es dedicarme a la física computacional cuántica. La aplicación que me gustaría aportar a la industria es mejorar la inteligencia artificial cuántica y/o crear algoritmos de computación cuántica para el pilotaje autónomo de aeronaves o estructuras robóticas.}

\end{document}
